% !TeX program = xelatex
\documentclass[UTF8,a4paper,zihao=-4,fontset=none]{ctexart}
\setCJKmainfont[BoldFont={LXGW Neo XiHei}]{LXGW Neo ZhiSong}
\setCJKsansfont[AutoFakeBold=2]{LXGW Neo XiHei}
\setCJKmonofont{LXGW Neo XiHei}
\setCJKfamilyfont{zhhei}[AutoFakeBold=2]{LXGW Neo XiHei}
\providecommand{\heiti}{\CJKfamily{zhhei}}
% Shared renderer entrypoint for generated lessonplan.tex files.
% Generated documents should load ctexart, then input this file.

\usepackage{hepingjie_lessonplan}
% Hepingjie lesson-plan overlay coordinates.
% Coordinate origin: top-left corner of an A4 portrait page.
% Units: mm. These coordinates are frozen for templates/lessonplan/hepingjie_blank.pdf.

\newcommand{\LPBackgroundPdf}{templates/lessonplan/hepingjie_blank.pdf}

% Page 1 fixed fields.
\newcommand{\LPDateYearX}{65mm}
\newcommand{\LPDateMonthX}{89mm}
\newcommand{\LPDateDayX}{110mm}
\newcommand{\LPDateY}{38mm}
\newcommand{\LPRecordPageX}{170mm}
\newcommand{\LPRecordPageY}{38mm}

\newcommand{\LPTopicX}{32mm}
\newcommand{\LPTopicY}{49mm}
\newcommand{\LPTopicW}{108mm}
\newcommand{\LPTypeX}{154mm}
\newcommand{\LPTypeY}{49mm}
\newcommand{\LPTypeW}{40mm}

\newcommand{\LPUnitTotalX}{67mm}
\newcommand{\LPUnitTotalY}{61mm}
\newcommand{\LPTopicTotalX}{111mm}
\newcommand{\LPTopicTotalY}{61mm}
\newcommand{\LPCurrentPeriodX}{169mm}
\newcommand{\LPCurrentPeriodY}{61mm}

\newcommand{\LPBodyX}{33mm}
\newcommand{\LPObjectiveY}{74mm}
\newcommand{\LPObjectiveH}{39mm}
\newcommand{\LPKeyPointY}{119mm}
\newcommand{\LPKeyPointH}{15mm}
\newcommand{\LPDifficultyY}{137mm}
\newcommand{\LPDifficultyH}{15mm}
\newcommand{\LPMethodY}{156mm}
\newcommand{\LPMethodH}{15mm}
\newcommand{\LPMediaY}{174mm}
\newcommand{\LPMediaH}{15mm}
\newcommand{\LPBlackboardY}{196mm}
\newcommand{\LPBlackboardH}{73mm}
\newcommand{\LPBodyW}{164mm}

% Teaching process boxes on pages 2-4.
\newcommand{\LPProcessX}{33mm}
\newcommand{\LPProcessY}{31mm}
\newcommand{\LPProcessW}{163mm}
\newcommand{\LPProcessFullH}{249mm}
\newcommand{\LPProcessPageFourH}{160mm}


\usetikzlibrary{calc}
\begin{document}

\LPBackgroundPage{1}{
  \LPCenterBox{\LPDateYearX}{\LPDateY}{12mm}{5mm}{}
  \LPCenterBox{\LPDateMonthX}{\LPDateY}{10mm}{5mm}{}
  \LPCenterBox{\LPDateDayX}{\LPDateY}{10mm}{5mm}{}
  \LPCenterBox{\LPRecordPageX}{\LPRecordPageY}{10mm}{5mm}{1}
  \LPHeaderCell{\LPTopicX}{\LPTopicY}{\LPTopicW}{8mm}{线段的比较与基本事实}
  \LPHeaderCell{\LPTypeX}{\LPTypeY}{\LPTypeW}{8mm}{几何探究课}
  \LPCenterBox{\LPUnitTotalX}{\LPUnitTotalY}{12mm}{5mm}{11}
  \LPCenterBox{\LPTopicTotalX}{\LPTopicTotalY}{12mm}{5mm}{3}
  \LPCenterBox{\LPCurrentPeriodX}{\LPCurrentPeriodY}{12mm}{5mm}{2}
  \LPTextBox{\LPBodyX}{\LPObjectiveY}{\LPBodyW}{\LPObjectiveH}{\LPPlainLine{说出并解释画一条线段等于已知线段，使用本节规范术语表达。}\LPPlainLine{按规范步骤完成线段比较，并说明关键依据。}\LPPlainLine{判断并订正与“尺规作图与线段比较方法的理解”有关的常见错误。}\LPPlainLine{在变式或实际情境中运用本节方法，完整写出过程和结论。}}
  \LPTextBox{\LPBodyX}{\LPKeyPointY}{\LPBodyW}{\LPKeyPointH}{线段比较；两点之间线段最短}
  \LPTextBox{\LPBodyX}{\LPDifficultyY}{\LPBodyW}{\LPDifficultyH}{尺规作图与线段比较方法的理解}
  \LPTextBox{\LPBodyX}{\LPMethodY}{\LPBodyW}{\LPMethodH}{观察操作；画图标注；图形说理}
  \LPTextBox{\LPBodyX}{\LPMediaY}{\LPBodyW}{\LPMediaH}{PPT；黑板；反馈卡；直尺；量角器或透明纸}
  \LPTextBox{\LPBodyX}{\LPBlackboardY}{\LPBodyW}{\LPBlackboardH}{\LPBoardTitle{线段的比较与基本事实}\LPBoardLine{1}{核心：可用叠合法或度量法比较；两点之间线段最短，线段长度叫两点间距离。}\LPBoardLine{2}{方法：叠合或度量比较，用线段最短解释路径问题。}\LPBoardLine{3}{例题：已知 $AB=5$ cm、$CD=3.5$ cm，比较。 答：$AB>CD$。}\LPBoardLine{4}{易错：判断：图上看起来更长的线段一定更长。 纠正：错误；图形仅示意，应依据度量、叠合或已知条件。}\LPBoardLine{5}{检查：条件、依据、步骤、符号或单位逐项核对。}}
}

\LPBackgroundPage{2}{
  \LPProcessBox{\LPProcessX}{\LPProcessY}{\LPProcessW}{\LPProcessFullH}{\LPProcessRow{观察与操作}{\LPTeacherPart{问题}{两条线段不便移动时，怎样比较长短？}\LPTeacherPart{组织}{出示题目或情境，先让学生独立判断，再请两名学生说明依据。}\LPTeacherPart{说明}{可用刻度尺度量，也可用圆规截取进行叠合比较。}\LPTeacherPart{纠错}{若学生只报结论，追问基准、条件或运算依据，要求补成完整句。}\LPTeacherPart{归纳}{把情境中的信息转化为数学对象和关系。}\LPTeacherPart{意图}{用具体问题激活先修经验，并明确本节需要解决的核心矛盾。}}{\LPStudentPart{活动}{独立观察或计算，圈出关键信息后口头说明。}\LPStudentPart{预期}{可用刻度尺度量，也可用圆规截取进行叠合比较。}}{5}
\LPProcessRow{概念形成}{\LPTeacherPart{问题}{线段比较和两点之间的基本事实是什么？}\LPTeacherPart{组织}{组织观察、比较或试算，板书学生给出的共同点，并逐项核对术语。}\LPTeacherPart{说明}{可用叠合法或度量法比较；两点之间线段最短，线段长度叫两点间距离。}\LPFigureBlock{
\begin{tikzpicture}[x=12mm,y=9mm,baseline=(current bounding box.center)]
  \draw (0,1.2)--(4,1.2); \fill(0,1.2)circle(1.2pt)node[above]{$A$};\fill(4,1.2)circle(1.2pt)node[above]{$B$};
  \draw (.5,0)--(3.2,0); \fill(.5,0)circle(1.2pt)node[above]{$C$};\fill(3.2,0)circle(1.2pt)node[above]{$D$};
  \node[right] at (4.5,.6){\scriptsize 叠合或度量比较};
\end{tikzpicture}
}\LPTeacherPart{例题}{基础例题：已知 $AB=5$ cm、$CD=3.5$ cm，比较。 规范解答：$AB>CD$。}\LPTeacherPart{对应练习}{即时练习：用圆规比较两条未标长度的线段。 答案：截取一条长度到另一条射线上观察端点位置。}\LPTeacherPart{纠错}{用反例检查是否真正满足条件，提醒按“叠合或度量比较，用线段最短解释路径问题”说明。}\LPTeacherPart{归纳}{叠合或度量比较，用线段最短解释路径问题}\LPTeacherPart{意图}{让核心结论由可检验的观察或运算形成，而不是直接记忆。}}{\LPStudentPart{活动}{在图形、算式或表格中标注条件，与同桌归纳共同特征。}\LPStudentPart{预期}{可用叠合法或度量法比较；两点之间线段最短，线段长度叫两点间距离。}}{7}
\LPProcessRow{图形表示}{\LPTeacherPart{问题}{解答“点 $A,B,C$ 不共线，比较 $AB+BC$ 与 $AC$。”前应先确定什么？}\LPTeacherPart{组织}{教师按题意、依据、步骤和结论顺序板演；学生在每一步旁标注作用。}\LPTeacherPart{说明}{规范解答：由两点之间线段最短，$AB+BC>AC$。}\LPTeacherPart{例题}{变式例题：点 $A,B,C$ 不共线，比较 $AB+BC$ 与 $AC$。 解答：由两点之间线段最短，$AB+BC>AC$。}\LPTeacherPart{对应练习}{对应练习：若 $AB=CD$，可用什么符号表示？ 答案：$AB=CD$。}\LPTeacherPart{纠错}{若一步跳到结论，要求补写关键关系或中间步骤，并用原题条件检验。}\LPTeacherPart{归纳}{解题流程：叠合或度量比较，用线段最短解释路径问题。}\LPTeacherPart{意图}{用完整例题建立可模仿的书写顺序和检查方法。}}{\LPStudentPart{活动}{先独立完成，再与板演对照步骤、符号、单位或图形标记。}\LPStudentPart{预期}{能得到 由两点之间线段最短，$AB+BC>AC$。，并说出关键依据。}}{7}}
}

\LPBackgroundPage{3}{
  \LPProcessBox{\LPProcessX}{\LPProcessY}{\LPProcessW}{\LPProcessFullH}{\LPProcessRow{关系应用}{\LPTeacherPart{问题}{从 $A$ 到 $B$ 有折线路径 $AC+CB=8$ km，直线路径 $AB=6$ km，选哪条更短？}\LPTeacherPart{组织}{要求学生先写条件关系或示意，再完成计算并解释结果；抽取两种方法比较。}\LPTeacherPart{说明}{直线路径 $AB$ 更短。}\LPTeacherPart{例题}{应用例题：从 $A$ 到 $B$ 有折线路径 $AC+CB=8$ km，直线路径 $AB=6$ km，选哪条更短？ 解答：直线路径 $AB$ 更短。}\LPTeacherPart{对应练习}{即时变式：从村庄到公路修最短道路体现什么事实？ 答案：两点之间线段最短。}\LPTeacherPart{纠错}{若答案脱离条件或单位，要求回到题干逐项对应并作合理性检查。}\LPTeacherPart{归纳}{先识别结构，再选择方法，最后解释结果。}\LPTeacherPart{意图}{把核心方法迁移到不同表示方式或实际情境中。}}{\LPStudentPart{活动}{独立列式或作图，小组内互查条件、步骤和答语。}\LPStudentPart{预期}{直线路径 $AB$ 更短。}}{6}
\LPProcessRow{条件辨析}{\LPTeacherPart{问题}{判断：图上看起来更长的线段一定更长。}\LPTeacherPart{组织}{展示错解，要求学生只圈第一处错误，写出依据后再完整订正。}\LPTeacherPart{说明}{错误；图形仅示意，应依据度量、叠合或已知条件。}\LPTeacherPart{例题}{易错辨析：判断：图上看起来更长的线段一定更长。}\LPTeacherPart{对应练习}{纠错要求：写出第一处错误，并按“叠合或度量比较，用线段最短解释路径问题”重做。参考：错误；图形仅示意，应依据度量、叠合或已知条件。}\LPTeacherPart{纠错}{不接受只写“错”或只改答案；必须指出错误对象、正确规则和订正过程。}\LPTeacherPart{归纳}{纠错顺序：定位错误、说明依据、规范订正、回题检验。}\LPTeacherPart{意图}{利用典型错误暴露概念边界、符号习惯或条件遗漏。}}{\LPStudentPart{活动}{独立找错、写依据、完成订正，再与同桌互评理由是否完整。}\LPStudentPart{预期}{错误；图形仅示意，应依据度量、叠合或已知条件。}}{5}
\LPProcessRow{综合探究}{\LPTeacherPart{问题}{三道练习分别考查哪个要点，先做哪一道能最快暴露问题？}\LPTeacherPart{组织}{组织限时题组：①用圆规比较两条未标长度的线段。 ②若 $AB=CD$，可用什么符号表示？ ③从村庄到公路修最短道路体现什么事实？；完成后按答案自查。}\LPTeacherPart{说明}{答案：①截取一条长度到另一条射线上观察端点位置。 ②$AB=CD$。 ③两点之间线段最短。}\LPTeacherPart{对应练习}{机动题：已知 $AB=5$ cm、$CD=3.5$ cm，比较。 学有余力者换一种方法说明；答案要点：$AB>CD$。}\LPTeacherPart{纠错}{正确率低于 70\% 时暂停机动题，按核心方法示范一道，再让学生订正同类错题。}\LPTeacherPart{归纳}{叠合或度量比较，用线段最短解释路径问题}\LPTeacherPart{意图}{集中检查方法稳定性，并为反馈测试前的即时补救提供依据。}}{\LPStudentPart{活动}{限时完成三题，按概念、步骤或应用给错因分类，并订正一题。}\LPStudentPart{预期}{能够依据“叠合或度量比较，用线段最短解释路径问题”完成题组并说清错因。}}{5}}
}

\LPBackgroundPage{4}{
  \LPProcessBox{\LPProcessX}{\LPProcessY}{\LPProcessW}{\LPProcessPageFourH}{\LPProcessRow{当堂检测}{\LPTeacherPart{问题}{四道反馈题中，哪一题最能检验本节核心方法？}\LPTeacherPart{组织}{投放 10 分反馈卡：①2 分，线段比较和两点之间的基本事实是什么？ ②2 分，若 $AB=CD$，可用什么符号表示？ ③2 分，判断：图上看起来更长的线段一定更长。 ④4 分，从 $A$ 到 $B$ 有折线路径 $AC+CB=8$ km，直线路径 $AB=6$ km，选哪条更短？}\LPTeacherPart{说明}{参考答案：①可用叠合法或度量法比较；两点之间线段最短，线段长度叫两点间距离。 ②$AB=CD$。 ③错误；图形仅示意，应依据度量、叠合或已知条件。 ④直线路径 $AB$ 更短。}\LPTeacherPart{纠错}{公布答案后学生自评并举牌统计：9 分及以上完成提高任务；7—8 分订正；7 分以下接受针对性补练。}\LPTeacherPart{归纳}{达标标准：10 分制 9 分及以上完全达标，7—8 分基本达标，7 分以下补练。}\LPTeacherPart{意图}{用概念、技能、纠错和综合四类任务覆盖本节目标。}}{\LPStudentPart{活动}{独立完成，按答案自评；把错误标为概念、步骤、条件或表达问题。}\LPStudentPart{预期}{能完成四类题并用核心术语说明至少一道题的依据。}}{6}
\LPProcessRow{回顾总结与作业}{\LPTeacherPart{问题}{本节研究了什么、关键步骤是什么、最容易错在哪里？}\LPTeacherPart{组织}{请学生各用一句话回答三个问题，教师据此补全板书并布置分层作业。}\LPTeacherPart{说明}{A 组：①用圆规比较两条未标长度的线段。 ②若 $AB=CD$，可用什么符号表示？ ③从村庄到公路修最短道路体现什么事实？ ④判断：图上看起来更长的线段一定更长。；答案：①截取一条长度到另一条射线上观察端点位置。 ②$AB=CD$。 ③两点之间线段最短。 ④错误；图形仅示意，应依据度量、叠合或已知条件。。B 组：①点 $A,B,C$ 不共线，比较 $AB+BC$ 与 $AC$。 ②从 $A$ 到 $B$ 有折线路径 $AC+CB=8$ km，直线路径 $AB=6$ km，选哪条更短？；答案要点：①由两点之间线段最短，$AB+BC>AC$。 ②直线路径 $AB$ 更短。。C 组选做：围绕“线段的比较与基本事实”自编一道含易错条件的题并给出解答与检查方法。}\LPTeacherPart{纠错}{作业答案必须包含必要步骤或依据；只写结果的题次日先口述方法再补写。}\LPTeacherPart{归纳}{核心方法：叠合或度量比较，用线段最短解释路径问题。课后反思区域由任课教师课后填写。}\LPTeacherPart{意图}{由学生完成结构化回顾，把课堂方法迁移到分层课后任务。}}{\LPStudentPart{活动}{说出一个结论、一个步骤和一个易错点，记录 A、B、C 三组作业。}\LPStudentPart{预期}{能用“叠合或度量比较，用线段最短解释路径问题”概括本节方法，并知道如何自查。}}{4}}
}
\end{document}
